
Phase 1: Identifying the "Collaboration Footprint"

Identify Heterogeneous Repositories: Select >100 open-source repositories that involve Agentic AI systems multi-agent systems, LLM pipelines (e.g., CrewAI, LangGraph, or ROS-based agent projects).
Artifact Harvesting: Extract not just the source code, but also the "collaboration footprint":
Orchestration Logic: Classes or functions that handle task delegation (e.g., Supervisor, Handoff, Process).
Communication Channels: Where data is passed between agents (e.g., MessageQueue, SharedState, ContextStore).
Validation Gates: Code that checks an agent's output before the next agent takes over.

Phase 2: Detecting the Smells (The Mining Task)

You should scan the repository for these specific smells that indicate a failing agentic architecture:
Some examples:

Erroneous data propagates through the pipeline, causing "compounding errors".
Multiple agents performing "Direct Mutation" on a shared state without locks.
Inconsistent assumptions and race conditions.	
Missing observability or logging traces for tool calls and reasoning.
Impossible to debug why an agent made a specific physical decision.	
Circular dependencies in agent planning.
Agentic systems getting stuck in a reasoning loop, draining resources.
Hard-coded prompts that lack explicit, measurable KPIs or goal models.
The agent's actions slowly diverge from the human's intended physical goal.


Phase 3: Agentic AI Approach for Smell Detection

Configure the "Detector" Agent: Use an Agentic AI framework to automate the search for "indicators".
